% Pursuing Stacks (cotes 134-2 to 134-8): archive cote 134-5, site page 72 (PDF page 73), typescript page 259.
% typescript. Draft compared: ps4.tex:3605-3692.
% AI-transcribed by Claude Opus 5.5 (claude-opus-5-5, adaptive thinking, effort high), Message Batches API, batch ps2 (msgbatch_01NVwr1Yuo8qJf5rykeGBXPK), 2026-11-12; reviewed against both supplied scan halves; see CHECK_DECISIONS.tsv.
\scanpage{134-5/72}
\mspage{259}

25.6.

81. I did'nt find much time since monday for mathematical pondering -
the little I got nonetheless has been enough though for convincing myself
that things came out \xout{\unsure{st}} more nicely still than I expected by then. \hand{\unsure{--}} \over{O}{\unsure{o}}ne
main point \(\underset{\hand{\cdot\ \cdot\ \cdot\ \cdot\ \cdot}}{\text{\deleted{being}}}\) that, provided the basic localizer \inserted[typed]{$\underline{W}$} satisfies the mild
extra assumption Loc 4) below, any contractibility structure $M_c$ on a
category $M$ \inserted[typed]{with finite products,} \xout{is} can be recovered, in the simplest imaginable way, in
terms of the associated asphericity structure $M_a$ , or equivalently, in
terms of the corresponding set $W_a$ of ``weak equivalences'';\xout{x} namely, $M_c$
is the set of contractible objects of $M$ , for the \inserted[typed]{homotopy} \xout{weak} interval structu-%
re defined in terms of all intervals made up with objects of $M_a$ \xout{x} . This
implies that the canonical map
\[ \hand{(1)} \qquad \mathrm{Homtp}_4(M) = \mathrm{Cont}(M) \hand{\hookrightarrow} \underline{W}\text{-}\mathrm{Asph}(M) \quad , \]
from the set of contractibility structures on $M$ to the set of aspherici-%
ty structures on $M$ relative to the basic localizer $\underline{W}$ (or \hand{``}\hand{\uline{\textcolor{black}{$\underline{W}$-asphericity
structures}}}\hand{''}), is \hand{\uline{\textcolor{black}{injective}}}. In other words, we may view \xout{thexnotionxof}
a contractibility structure (on a category $M$ stable under finite products\hand{),}
which is an \hand{\uline{\textcolor{black}{absolute}}} notion (namely independaht of \xout{any}\inserted{the} choice of \xout{the bax}
\inserted{a} basic localizer $\underline{W}$), as a``\hand{\uline{\textcolor{black}{particular}}} \hand{\uline{\textcolor{black}{case}}}'' of a $\underline{W}$-asphericity structure
(depending on the choice of $\underline{W}$), namely as ``equivalent'' to a $\underline{W}$-aspherici-%
ty structure, satisfying some extra conditions which we'll have to writ\hand{e}
down below. \hand{//}
% The paragraph-break mark has two near-vertical strokes with horizontal extensions beneath down below and above the following This; the upper extension lies near satisfying.

This pleasant fact associates immediately with two related
ones. The first is just a reminder of our reflections of sections 51 and
52, namely that the set of contractibility structures on $M$ can be viewed
as one among four similar sets of ``homotopy structures'' on $M$
\[ \hand{(2)} \qquad \mathrm{Homtp}_4(M) \hand{\hookrightarrow} \mathrm{Homtp}_3(M) \hand{\hookrightarrow} \mathrm{Homtp}_2(M) \hand{\xrightarrow{\sim}} \mathrm{Homtp}_1(M) \quad , \]
corresponding to the four \xout{maxim} basic ``homotopy notions'' met with so far,
namely (besides contractibility structures $\mathrm{Homtp}_4$) the \xout{weak} \hand{\uline{\textcolor{black}{homotopy
interval structures}}} ($\mathrm{Homtp}_3$), the \hand{\uline{\textcolor{black}{homotopism structures}}} \xout{$W_o$} ($\mathrm{Homtp}_2$) and
the \hand{\uline{\textcolor{black}{homotopy relations}}} \inserted{between maps} ($\mathrm{Homtp}_1$) . In the sequel, if \xout{$M_c$xxxObxMxxis}
\[ \hand{(3)} \qquad M_c \hand{\subset} \mathrm{Ob}\, M \]
is a given homotopy structure on $M$, we'll denote by \xout{\unsure{$J_c$}}
\[ \hand{(4)} \qquad J_c \hand{\subset} \mathrm{Int}(M), \quad W_c \hand{\subset} \mathrm{Fl}(M) \ , \quad R_c \hand{\subset} \mathrm{Fl}(M)\text{x}\mathrm{Fl}(M) \]
the corresponding other three homotopy structures on $M$, where $\mathrm{Int}(M)$
denotes the set of all``intervals'' $\underline{I} = (I, \hand{\delta}_o, \hand{\delta}_1)$ in $M$, i.e. objects of
$M$ endowed with two sections $\hand{\delta}_o,\hand{\delta}_1$ over the final object $e_M$ of $M$.

\margin{72}
